%% Chapter 8 — Section 8.2: From Hamiltonian Dynamics to a Sampling Algorithm
%% Source: PDF pages 105–107 (book pages 92–94)

\section{From Hamiltonian Dynamics to a Sampling Algorithm}
\label{sec:8.2}

%% -------------------------------------------------------
\subsection{Idealized Hamiltonian Monte Carlo Algorithm}
\label{sec:8.2.1}
%% -------------------------------------------------------

For the following algorithm, we consider a Hamiltonian of the form
$H(q,p) = \frac{1}{2}p^T M^{-1}p + V(q)$. We would like to construct a Markov kernel that
leaves the canonical measure of $H$ invariant. Recall that our true target is the $q$-marginal
with density proportional to $e^{-V}$; the Gaussian momentum $p$ is just an auxiliary variable.
As usual, we have access to the potential $V(q)$ (and its gradient) but not to the normalizing
constant $Z$.

\begin{algorithm}[H]
\caption{Idealized, Not-Implementable Hamiltonian Monte Carlo}
\label{alg:8.1}
\begin{algorithmic}[1]
\Require Initialization $(q^{(0)}, p^{(0)})$, duration parameter $\lambda$, sample size $N$.
\For{$n = 0, \ldots, N-1$}
  \State \textbf{Momentum refreshment:} Sample $p^* \sim \mathcal{N}(0, M)$.
  \State \textbf{State update:} $(q^{(n+1)}, p^{(n+1)}) = \varphi_\lambda(q^{(n)}, p^*)$.
\EndFor
\Ensure Sample $\{q^{(n)}\}_{n=1}^N$.
\end{algorithmic}
\end{algorithm}

Note that the idealized, not implementable HMC algorithm shown above presupposes that the flow
map $\varphi_\lambda$ can be evaluated, which is not true for most practical applications.
However, it is important to note that if the flow map could be evaluated, there would be no need
for an acceptance/rejection step to leave the target invariant; this is proved in
Theorem~\ref{thm:8.10}. In practice, $\varphi_\lambda$ is approximated with numerical
integrators (which might not preserve total energy) and the acceptance/rejection step is used to
counter that error. This will be made precise in Algorithm~\ref{alg:8.2} and
Theorem~\ref{thm:8.12}.

\begin{theorem}
\label{thm:8.10}
Let $\mu^H$ be the Gibbs distribution with energy $H(q,p) = \frac{1}{2}p^T M^{-1}p + V(q)$.
The Markov kernel implicitly defined by Algorithm~\ref{alg:8.1} leaves $\mu^H$ invariant.
\end{theorem}

\begin{proof}
By the definition of the Hamiltonian, the momentum variables are $\mathcal{N}(0,M)$ and
therefore $\mu^H$ is clearly invariant in the momentum refreshment step. On the other hand,
Theorem~8.5 shows that the flow map $\varphi_\lambda$ is measure-preserving with respect to
the canonical measure of $H$. This implies that step~2 also preserves the distribution $\mu^H$.
\end{proof}

%% -------------------------------------------------------
\subsection{Hamiltonian Monte Carlo Algorithm}
\label{sec:8.2.2}
%% -------------------------------------------------------

Time reversibility, conservation of Hamiltonian, and symplecticity are properties that hold for
Hamiltonian systems. In practice, for most problems of interest, Hamilton equations cannot be
solved in closed form and need to be discretized. Discretizations of Hamiltonian systems may
preserve \emph{exactly} some of these properties. For instance, the famous leapfrog method is
reversible and symplectic. We remark that unfortunately a numerical solver for Hamilton
equations cannot be simultaneously symplectic \emph{and} conserve the Hamiltonian exactly.

\begin{example}[The Leapfrog Integrator]
\label{ex:8.11}
The leapfrog method applied to Hamilton equations with
$K(p) := \sum_{i=1}^d \frac{p_i^2}{2m_i}$ and step-size $\epsilon > 0$ is
\begin{align*}
p_i\!\left(t + \tfrac{\varepsilon}{2}\right)
  &= p_i(t) - \frac{\varepsilon}{2}\frac{\partial V}{\partial q_i}(q(t))
  &\quad&\text{(half momentum step)} \\[4pt]
q_i(t + \varepsilon)
  &= q_i(t) + \varepsilon \frac{p_i\!\left(t + \frac{\varepsilon}{2}\right)}{m_i}
  &\quad&\text{(full position step)} \\[4pt]
p_i(t + \varepsilon)
  &= p_i\!\left(t + \tfrac{\varepsilon}{2}\right)
    - \frac{\varepsilon}{2}\frac{\partial V}{\partial q_i}(q(t+\varepsilon)).
  &\quad&\text{(half momentum step)}
\end{align*}
It is time reversible and conserves volume exactly. To discretize Hamilton equations for a time
interval $[t, t+\lambda]$ of length $\lambda$, we may take $L$ leapfrog steps with
$\lambda = \varepsilon L$.
\end{example}

\begin{figure}[htbp]
\centering
\includegraphics[width=0.95\textwidth]{figures/fig_08_2}
\caption{\textit{Simulate the Hamiltonian $\frac{3p^2}{2} + \frac{4q^2}{2}$ using leapfrog
and Euler method updates: $p(t+\varepsilon) = p(t) + \varepsilon\frac{dp}{dt}$,
$q(t+\varepsilon) = q(t) + \varepsilon\frac{dq}{dt}$.}}
\label{fig:8.2}
\end{figure}

Let $\Psi_\lambda$ be a symplectic numerical integrator, time reversible with respect to the
momentum flip involution, that approximates the flow map $\varphi_\lambda$. The Hamiltonian
Monte Carlo algorithm proposes moves using $\Psi_\lambda$ and corrects for the fact that
$\Psi_\lambda$ does not preserve the Hamiltonian by using an accept/reject mechanism.

\begin{algorithm}[H]
\caption{Hamiltonian Monte Carlo (HMC)}
\label{alg:8.2}
\begin{algorithmic}[1]
\Require Initialization $(q^{(0)}, p^{(0)})$, duration parameter $\lambda$, sample size $N$.
\For{$n = 0, \ldots, N-1$}
  \State \textbf{Momentum refreshment:} Sample $p^{\mathrm{new}} \sim \mathcal{N}(0, M)$.
  \State \textbf{Propose new state:} $(q^*, p^*) = \Psi_\lambda(q^{(n)}, p^{\mathrm{new}})$.
  \State \textbf{Set}
  \[
    (q^{(n+1)}, p^{(n+1)}) =
    \begin{cases}
      (q^*, p^*) & \text{with probability }
        a := \min\!\left\{1,\, e^{H(q^{(n)}, p^{\mathrm{new}}) - H(q^*, p^*)}\right\}, \\
      (q^{(n)}, -p^{\mathrm{new}}) & \text{with probability } 1 - a.
    \end{cases}
  \]
\EndFor
\Ensure Sample $\{q^{(n)}\}_{n=1}^N$.
\end{algorithmic}
\end{algorithm}

\begin{center}
\fbox{\parbox{0.92\textwidth}{%
\textbf{Pros and Cons}: Using Hamiltonian dynamics in an extended state space, HMC breaks away
from the local behavior of RWMH and MALA proposals. As a result, HMC scales better to
high-dimensional problems (see below for further discussion). However, HMC is often harder to
implement and to tune than RWMH and MALA. Similar to MALA, HMC requires evaluation of the
gradient of $V$, which can be expensive.}}
\end{center}

\medskip
\noindent\textbf{Scaling of HMC} \quad
Classical theoretical results for HMC developed in the same large-$d$ scaling setting that we
considered for RWMH in Chapter~4 and for MALA in Chapter~6 revealed the scaling of the
leapfrog step-size should be $\mathcal{O}(d^{-1/4})$. Thus, HMC requires $\mathcal{O}(d^{1/4})$
steps to traverse the state space, which should be compared to $\mathcal{O}(d^{1/3})$ for MALA
and $\mathcal{O}(d^1)$ for RWMH. These analyses also show that the optimal acceptance rate for
HMC is $0.651$, compared to $0.574$ for MALA and $0.234$ for RWMH.

Our goal in the rest of this section is to understand why the acceptance probability in the HMC
algorithm is the correct one to ensure that $f^H$ invariant is an invariant distribution. To
that end, we will make use of the time reversibility of the numerical integrator and its
preservation of volume (otherwise a Jacobian would appear in the acceptance probability). The
main theoretical result of this chapter is the following.

\begin{theorem}[HMC Leaves Target Invariant]
\label{thm:8.12}
Suppose that $\Psi_\lambda$ is symplectic and reversible with respect to the momentum flip
involution. Then, the Markov kernel implicitly defined by HMC leaves invariant the canonical
measure of $H$.
\end{theorem}

\begin{proof}
Throughout the proof, we denote by $S$ the momentum flip involution. First, notice that the
momentum refreshment step preserves the canonical measure $\mu^H$. Hence, it suffices to show
that $\mu^H$ is invariant under steps~2 and~3. The corresponding Markov kernel is given by
\[
\pi_x(dy) = a(x)\,\delta_{\Psi_\lambda(x)}(y)\,dy
+ \bigl(1 - a(x)\bigr)\,\delta_{S(x)}(y)\,dy,
\]
where
\[
a(x) := \min\!\left\{1,\, e^{H(x) - H(\Psi_\lambda(x))}\right\}.
\]
Let $A \subset \R^{2d}$ be a Lebesgue measurable set. We need to show that
\[
\int_{\R^{2d}} \pi_x(A)\,\mu^H(dx) = \int_A \mu^H(dx).
\]
Expanding the left-hand side gives
\begin{align*}
\int_{\R^{2d}} \pi_x(A)\,\mu^H(dx)
&= \int_{\R^{2d}} a(x)\,\mathbf{1}_A\!\left(\Psi_\lambda(x)\right) f^H(x)\,dx
  + \int_{\R^{2d}} \bigl(1-a(x)\bigr)\,\mathbf{1}_A\!\left(S(x)\right) f^H(x)\,dx \\
&= \int_{\R^{2d}} \mathbf{1}_A\!\left(S(x)\right) f^H(x)\,dx
  + \int_{\R^{2d}} a(x)\,\mathbf{1}_A\!\left(\Psi_\lambda(x)\right) f^H(x)\,dx
  - \int_{\R^{2d}} a(x)\,\mathbf{1}_A\!\left(S(x)\right) f^H(x)\,dx.
\end{align*}
First, $\int_{\R^{2d}} \mathbf{1}_A(S(x)) f^H(x)\,dx = \int_A \mu^H(dx)$ since $S$ preserves
the canonical measure. Next, we use a change of variables in the third term to show that it is
equal to the second one. Precisely, using that $S \circ S = \mathrm{id}$, and that the
determinant Jacobians of $\Psi_\lambda$ and $S$ are identically one due to the symplecticity of
$\Psi_\lambda$ and reversibility of $S$, we have
\[
\int_{\R^{2d}} a(x)\,\mathbf{1}_A\!\left(S(x)\right) f^H(x)\,dx
= \int_{\R^{2d}} a\!\left(S \circ \Psi_\lambda(x)\right)
  \mathbf{1}_A\!\left(\Psi_\lambda(x)\right) f^H\!\left(\Psi_\lambda(x)\right) dx.
\]
Now, using that $\Psi_\lambda \circ S \circ \Psi_\lambda = S$ since $\Psi_\lambda$ is time
reversible with respect to $S$ and also that $H \circ S = H$, we deduce that
\begin{align*}
a\!\left(S \circ \Psi_\lambda(x)\right)
&= \min\!\left\{1,\, e^{H(S \circ \Psi_\lambda(x)) - H(\Psi_\lambda \circ S \circ \Psi_\lambda(x))}\right\} \\
&= \min\!\left\{1,\, e^{H(\Psi_\lambda(x)) - H(S(x))}\right\} \\
&= \min\!\left\{1,\, e^{H(\Psi_\lambda(x)) - H(x)}\right\} \\
&= \frac{f^H(x)}{f^H\!\left(\Psi_\lambda(x)\right)}\cdot a(x).
\end{align*}
Substituting this back to the previous integral gives
\[
\int_{\R^{2d}} a(x)\,\mathbf{1}_A\!\left(S(x)\right) f^H(x)\,dx
= \int_{\R^{2d}} a(x)\,\mathbf{1}_A\!\left(\Psi_\lambda(x)\right) f^H(x)\,dx,
\]
as desired.
\end{proof}
